\section{Other notable programs}
There are two other pieces of software that also should be mentioned;
\pulledpork{} and \oinkmaster{} are both programs that can download rules and
update \snort{}-sensors. These programs also contain some of the features found
in \snowman{} and \bhtb{}, and \pulledpork{} even supports
Shared Object rules \cite{SnortManual}, which is not part of this project due to
complexity.  However, both of these programs suffer from the fact that they are
command-line oriented. They are both very good at what they do, but lack the
usability and efficiency offered by GUIs with buttons. These programs are a
good choice for rule management in a small environment with one or very few
sensors, but once the environment grows larger, with tens or even hundreds of
sensors, they become ineffective.

\section{Research on Software Performance}
Since this project also investigates software performance, it is relevant to
take a look at what others in the field write on the subject. This subject is
missing from most software development reports from \GUC{}, which can be
explained by the usual high focus on agile\cite{AGILE} development. 
\cite{FutureEngineering} mentions that agile programming is becoming mainstream
and by its nature is a concern for the future of software performance
engineering, as the primary focus of an agile project is on functionality.

A proper performance testing methodology was hence not immediately available for
this project. However, the articles \cite{WOSP98} and \cite{WOSP04} gave a few
pointers towards developing the test cases described in chapter
\ref{chap:performance_analysis}. Both articles also stress the importance of
performance testing during the development rather than after it, as this allows
for time to actually implement optimizations and improvements. The performance
testing in this report does not strictly follow any methodology, but is a
result of the team members' use of common sense and the scientific method.

Early on in the project it became clear that \snowman{} is very I/O bound, as it does
a lot of database operations. The tests in this report show a noticable
difference in performance between traditional rotational hard drives and SSDs.
Myers speaks of this in his Master thesis \cite{Myers} which shows that SSD
has a significant advantage over HDDs when it comes to random read performance
but not when it comes to write performance. Note that these tests were performed
in 2008, a time where SSD technology was fairly new and suffering
from their different way of writing data\footnote{Features such as TRIM was
introduced shortly after Myers thesis.}. As operating systems, drivers and disk
controllers has improved, SSDs are now superior to HDDs in most tasks.